% !TEX program = xelatex
\documentclass{英语试卷}
\begin{document}
\examtitle{2025 届湖北省高三(4 月)调研模拟考试}{英语}
\par 英语学科 试题卷\par
\par 命题单位：宜昌市教科院 审题单位：襄阳市教科院 黄冈市教科院\par
\begin{notice}
\par 1. 答卷前，考生务必将自己的姓名、准考证号填写在答题卡上。\par
\par 2. 回答选择题时，选出每小题答案后，用铅笔把答题卡上对应题目的答案标号涂黑。如需改动，用橡皮擦干净后，再选涂其他答案标号。回答非选择题时，将答案写在答题卡上，写在本试卷上无效。\par
\par 3. 考试结束后，将本试卷和答题卡一并交回。\par
\end{notice}
\parttitle{第一部分 听力（共两节，满分30分）}
\par 做题时，先将答案标在试卷上，录音内容结束后，你将有两分钟的时间将试卷上的答案转涂到答题卡上。\par
\sectitle{第一节（共5小题：每小题1.5分，满分7.5分）}
\par 听下面5段对话。每段对话后有一个小题，从题中所给的A、B、C三个选项中选出最佳选项。听完每段对话后，你都有10秒钟的时间来回答有关小题和阅读下一小题。每段对话仅读一遍。\par
\question{1. How does the woman feel?}
\begin{choices}(3)
  \item Confused.
  \item Stressed.
  \item Satisfied.
\end{choices}
\question{2. When will the new National Geographic probably arrive?}
\begin{choices}(3)
  \item On Wednesday.
  \item On Thursday.
  \item On Friday.
\end{choices}
\question{3. What is the woman complaining about?}
\begin{choices}(3)
  \item Her job.
  \item A report.
  \item A phone call.
\end{choices}
\question{4. Which place attracted the man most?}
\begin{choices}(3)
  \item The church.
  \item The art gallery.
  \item The botanical garden.
\end{choices}
\question{5. Why does the man come to the woman?}
\begin{choices}(3)
  \item To fill out a form.
  \item To borrow a book.
  \item To get a library card.
\end{choices}
\sectitle{第二节（共15小题：每小题1.5分，满分22.5分）}
\par 听下面5段对话或独白。每段对话或独白后有几个小题，从题中所给的A、B、C三个选项中选出最佳选项。听每段对话或独白前，你将有时间阅读各个小题，每小题5秒钟；听完后，各小题将给出5秒钟的作答时间。每段对话或独白读两遍。 听第6段材料，回答第6、7题。\par
\question{6. What is the possible relationship between the speakers?}
\begin{choices}(3)
  \item Friends.
  \item Coworkers.
  \item Strangers.
\end{choices}
\question{7. What kind of films does the woman like best?}
\begin{choices}(3)
  \item Documentaries.
  \item Action films.
  \item Sci-fi films.
\end{choices}
\par 听第7段材料，回答第8、9题。\par
\question{8. What is the woman's research paper about?}
\begin{choices}(3)
  \item Renewable energy.
  \item The southern region.
  \item Traditional architecture.
\end{choices}
\question{9. Where did the woman get materials for her paper?}
\begin{choices}(3)
  \item From books.
  \item From the Internet.
  \item From journals.
\end{choices}
\par 听第8段材料，回答第10至12题。\par
\question{10. What is the man's favourite chocolate?}
\begin{choices}(3)
  \item Milk chocolate.
  \item White chocolate.
  \item Dark chocolate.
\end{choices}
\question{11. Where does the woman believe the best chocolate comes from?}
\begin{choices}(3)
  \item France.
  \item Belgium.
  \item Germany.
\end{choices}
\question{12. What will the man do next?}
\begin{choices}(3)
  \item Go home.
  \item Go to a supermarket.
  \item Go to the chocolate workshop.
\end{choices}
\par 第9段材料，回答第13至16题。\par
\question{13. How often are private math lessons offered?}
\begin{choices}(3)
  \item Once a week.
  \item Twice a week.
  \item Twice a month.
\end{choices}
\question{14. What classes are given on Thursday evenings?}
\begin{choices}(3)
  \item Tennis classes.
  \item Math classes.
  \item Yoga classes.
\end{choices}
\question{15. What does the man say about the book club?}
\begin{choices}(3)
  \item It is a new club.
  \item It is aimed at teens.
  \item It is closed at weekends.
\end{choices}
\question{16. How many days per week is the centre open?}
\begin{choices}(3)
  \item Five days.
  \item Six days.
  \item Seven days.
\end{choices}
\par 第10段材料，回答第17至20题。\par
\question{17. What is new in Brown Hall?}
\begin{choices}(3)
  \item The gym.
  \item The library.
  \item The restaurant.
\end{choices}
\question{18. Which hall is accessible only to girls?}
\begin{choices}(3)
  \item Queens Building.
  \item Blake Residence.
  \item The Parkway Flats.
\end{choices}
\question{19. What do we know about the Parkway Flats?}
\begin{choices}(3)
  \item It has private showers.
  \item It is famous for its design.
  \item It is the most expensive hall.
\end{choices}
\question{20. What might attract students to Temple Rise?}
\begin{choices}(3)
  \item The cost of meals.
  \item The large rooms.
  \item The cleaning service.
\end{choices}
\parttitle{第二部分 阅读理解（共两节，满分50分）}
\sectitle{第一节（共15小题：每小题2.5分，满分37.5分）}
\par 阅读下列短文，从每题所给的A、B、C和D四个选项中，选出最佳答案。\par
\passagelabel{A}
\par Community Garden Volunteers Needed!\par
\par Are you passionate about sustainability and community building? Green Horizons, a local non-profit organization, is looking for enthusiastic volunteers to join our program!\par
\par What You'll Do:\par
\par Educate visitors about the benefits of locally grown food. Assist with planting, weeding, and harvesting fresh produce. Maintain garden tools and ensure the garden is a safe and welcoming space for all. Support the delivery of harvested produce to local food banks and community members in need.\par
\par What We're Looking For:\par
\par Commitment to making a difference in the community. Essential gardening know-how and willingness to learn. Ability to work both independently and as part of a team. A positive attitude and a passion for environmental sustainability.\par
\par What You'll Gain: Hands-on experience in sustainable gardening and community service.\par
\par A certificate of appreciation and a reference letter for your volunteer work. The satisfaction of contributing to a meaningful cause that benefits both people and the planet.\par
\par Commitment:\par
\par Volunteers are asked to commit at least 4 hours per week during the program period. Shifts can be tailored to fit your schedule.\par
\par How to Apply:\par
\par Send your application to volunteer@greenhorizons.org or visit our website for more details. Join Us Today and Help Cultivate a Better Tomorrow!\par
\question{21. What will the volunteers do?}
\begin{choices}(2)
  \item Maintain local food banks.
  \item Deliver locally grown produce.
  \item Teach sustainable gardening.
  \item Guide safe use of garden tools.
\end{choices}
\question{22. What is required of the volunteers?}
\begin{choices}(2)
  \item Willingness to share.
  \item Commitment to charity.
  \item Strong leadership skills.
  \item Basic gardening knowledge.
\end{choices}
\question{23. What do we know about the scheduling of volunteer shifts?}
\begin{choices}(2)
  \item It is regular.
  \item It is balanced.
  \item It is flexible.
  \item It is weekend-only.
\end{choices}
\passagelabel{B}
\par Nathan Sawaya is indeed known for his impressive artwork created using LEGO bricks. Initially, he faced legal challenges from LEGO regarding the use of its products in his art. However, over time, he gained recognition and became a "LEGO Certified Professional", which allows him to create and display his work with the company's support. This title is given to select individuals who demonstrate exceptional skill in working with LEGO and contribute to the company positively. Sawaya's unique approach to art has helped bridge the gap between commercial products and creative expression.\par
\par Sawaya's journey began in a rather unconventional way. Originally a corporate lawyer, he found himself increasingly drawn to the world of art. The turning point came when he decided to leave his stable job and pursue his passion for creating with LEGO. His first major exhibition, "The Art of the Brick", displayed his innovative sculptures and earned widespread praise. Visitors were amazed to see iconic artworks, such as Van Gogh's Starry Night and Michelangelo's David, reimagined entirely with LEGO bricks.\par
\par What sets Sawaya apart is not just his technical skill, but also his ability to express deep emotions through playful materials. Each piece tells a story, inviting viewers to engage with their inner child while thinking about the themes of creativity and imagination. His work challenges traditional concepts of art, demonstrating that even the most ordinary items can be turned into extraordinary expressions of human experience.\par
\par In addition to his exhibitions, Sawaya has also conducted workshops and educational programs, inspiring the next generation of artists. He believes that creativity should be accessible to everyone and encourages people to explore their artistic potential using simple materials. Through his devotion and passion, Nathan Sawaya has not only carved out a unique space in the art world but has also encouraged a greater appreciation for the creative possibilities found in LEGO.\par
\question{24. What does LEGO Certified Professional mean to Sawaya?}
\begin{choices}(2)
  \item It expands his name recognition.
  \item It legalizes his use of LEGO.
  \item It boosts his artistic competence.
  \item It finances his LEGO creation.
\end{choices}
\question{25. Which can best describe Sawaya's first major exhibition?}
\begin{choices}(2)
  \item Artistic and practical.
  \item Professional and vivid.
  \item Novel and impressive.
  \item Realistic and fashionable.
\end{choices}
\question{26. How does Sawaya's work challenge traditional views of art?}
\begin{choices}(2)
  \item By transforming ordinary objects into art.
  \item By featuring engagement and imagination.
  \item By taking full advantage of playful materials.
  \item By combining inner thought with technology.
\end{choices}
\question{27. What can we learn from Sawaya's example?}
\begin{choices}(2)
  \item Follow your passion, and success will follow.
  \item Share your knowledge, and wisdom will grow.
  \item Achievement is guaranteed for those with talent.
  \item Legal recognition is essential for artistic freedom.
\end{choices}
\passagelabel{C}
\par A battle broke out at the weekend over a hedge (树篱) in Lincolnshire. The hedge, near the town of Winterton, was covered in netting by Partner Construction, which has applied for planning permission to build 40 homes on the site. This is standard practice, the developer said, in order to prevent birds from nesting in a habitat that might be damaged if building work begins later in the year.\par
\par However, a group of local people opposed to the development produced a video showing birds trapped beneath the nets. Chris Packham shared the video on Twitter. Packham said the nets showed disregard of how to look after the countryside, and said, if he were there, he would "tear those nets".\par
\par The Wildlife and Countryside Act 1981 makes it illegal to intentionally "take, damage or destroy" the nest of any wild bird while it is in use or being built. In theory, this even bans people from cutting away some of the branches of their fruit trees if doing so harms a nest that they know is there. Any tree or hedge becomes a protected site the moment a bird settles in it, which is often the case between February and August—potentially shutting down an entire construction site for the busy half of the year.\par
\par To address the problem, conservationists and builders agree to find a middle ground. Netting an empty tree or a bush in winter is certainly better than making the birds look for a new site in the middle of the reproduction season. "Netting is probably the lowest-impact option if a habitat has to be removed during the reproduction season," a spokesman for Partner Construction said. "We appointed qualified and experienced ecologists to carry out the work on our behalf. They have convinced us and provided photographic evidence that there were no birds in the hedge after it was netted."\par
\question{28. What is the standard practice according to paragraph 1?}
\begin{choices}(2)
  \item Protecting birds with hedges.
  \item Applying for planning permission.
  \item Removing hedges from the site.
  \item Netting permitted construction sites.
\end{choices}
\question{29. What does the underlined word "disregard" in paragraph 2 mean?}
\begin{choices}(2)
  \item Evidence.
  \item Prospect.
  \item Ignorance.
  \item Inefficiency.
\end{choices}
\question{30. What do we know about the Wildlife and Countryside Act 1981?}
\begin{choices}(2)
  \item It protects wild birds' nests.
  \item It permits proper construction.
  \item It guides tree cutting practice.
  \item It prioritizes bird nesting seasons.
\end{choices}
\question{31. What does the last paragraph mainly talk about?}
\begin{choices}(2)
  \item The settlement of the conflict.
  \item The influence of the battle.
  \item The process of making a deal.
  \item The significance of netting a hedge.
\end{choices}
\passagelabel{D}
\par Surgeons often emphasize that in the event of a heart attack, time is crucial for saving heart muscle. The heart relies on a continuous supply of oxygen delivered by the heart's blood vessels (血管). When these vessels become blocked, the oxygen supply stops, and the heart's muscle cells begin to die within minutes. Survivors of heart attacks frequently suffer from permanent heart failure, with a worrying statistic showing that 50\% of them may not survive beyond five years. "Eventually, their hearts become so weak that they can’t sustain sufficient blood flow, and they ultimately stop," explains Sanjay Sinha, a cardiologist at Addenbrooke’s Hospital in Cambridge.\par
\par However, within the next five years, there could be a new alternative—to grow live, beating "heart patches" (补丁). Sinha, along with a team of biologists at the University of Cambridge’s Stem Cell Institute, is working on this innovative approach. These heart patches are tiny, beating pieces of heart muscle. They are grown in lab dishes over a month. The process begins by taking blood cells and reprogramming them into a specific type of stem cell capable of developing into any cell type within the human body. Sinha is currently preparing to test these patches, first in mice and then in pigs. If successful, he hopes to set up a human trial within five years.\par
\par In the United States, scientists from Stanford University are also working on creating heart patches. They envision a future procedure where ultrasound and MRI (超声波和核磁共振) scans are used to identify damaged sections. Based on the shape of the section, a custom heart patch would be 3D printed to fit perfectly. Surgeons would then open the chest area and attach the patch directly to the heart.\par
\par The major challenge with this method is establishing "electrical connectivity" between the new patch and the existing heart to ensure they beat in harmony. Additionally, the cost of developing and using these patches is significant, which could impact accessibility. However, if these challenges can be overcome, Sinha is optimistic that heart patches could transform the lives of millions worldwide.\par
\question{32. What is the function of paragraph 1?}
\begin{choices}(2)
  \item To reveal a finding.
  \item To offer a solution.
  \item To present a problem.
  \item To outline a research.
\end{choices}
\question{33. What is the first step in growing heart patches?}
\begin{choices}(2)
  \item Launching a clinical trial in humans.
  \item Testing the heart patches on animals.
  \item Culturing heart patches in lab dishes.
  \item Generating stem cells from blood cells.
\end{choices}
\question{34. What is the purpose of using ultrasound and MRI scans?}
\begin{choices}(2)
  \item To create a 3D model of the heart.
  \item To locate damaged areas in the heart.
  \item To design the shape of the custom heart.
  \item To provide real-time guidance for operations.
\end{choices}
\question{35. What is the major challenge in employing heart patches?}
\begin{choices}(2)
  \item Promoting global accessibility.
  \item Reducing production costs.
  \item Securing functional integration.
  \item Finding a suitable material.
\end{choices}
\sectitle{第二节（共5小题；每小题2.5分，满分12.5分）}
\par 阅读下面短文，从短文后的选项中选出能填入空白处的最佳选项。选项中有两项为多余选项。\par
\par As an enthusiastic reader, I've always believed that the key to understanding the world lies in the pages of a book. \blanknum{36} They introduce us to characters who feel as real as the people around us. This connection between readers and the text is what makes literature so powerful, and it's why I value my personal library. It's a collection that reflects my journey through the world of words.\par
\par My library is not vast, but each book is carefully chosen. I have a mix of types in my collection, ranging from classic literature to contemporary fiction. \blanknum{37} Each type offers a different perspective, a new way of looking at life. While some books challenge my thinking, others provide comfort. It's important to have a balance. A well-organized library is like a well-balanced diet for the mind.\par
\par Among my favourites are the works of Jane Austen and George Orwell. Austen's sharp observations on society and relationships are timeless. Her characters are imperfect yet charming. Orwell, on the other hand, offers a stark (冷峻的) view of the world. \blanknum{38} Both authors, though different in style, have taught me invaluable lessons\par
\par about humanity.\par
\par \blanknum{39} Biographies (传记) and life stories offer insights into the lives of extraordinary individuals. They remind us that greatness comes in many forms. Reading about the struggles and successes of others can be incredibly inspiring. It's a reminder that we all have the potential to make a difference. These stories encourage me to reflect on my own life and what I wish to leave behind.\par
\par My library is a reflection of who I am. It's a sign of my love of learning and my desire to understand the world. Each book is a stepping stone on my journey of discovery. They are companions in loneliness and guides in times of uncertainty. \blanknum{40}\par
\begin{choices}(1)
  \item I also enjoy non-fiction.
  \item I believe in the power of variety.
  \item My library goes beyond a collection of books.
  \item His stories are a warning and a call to action.
  \item Through them, I continue to explore, to grow, and to dream.
  \item The stories we read shape our world view and broaden our minds.
  \item Historical fiction, in particular, holds a special place in my heart.
\end{choices}
\parttitle{第三部分 语言知识运用（共两节，满分30分）}
\sectitle{第一节（共15小题；每小题1分，满分15分）}
\par 阅读下面短文，从短文后各题所给的A、B、C和D四个选项中，选出可以填入空白处的最佳选项。\par
\par I grew up in poverty, and that was pretty obvious to all my teachers. In 8 th grade, I won a contest for \blanknum{41} the yearbook cover, securing my \blanknum{42} as "that one art kid" in my class. I made the design with random pencils and crayons in the classroom, because I didn't have any of my own.\par
\par My teachers noticed I \blanknum{43} enjoyed making art, because there would always be doodles (涂鸦) in my notes. Everyone wanted to \blanknum{44} with me in any assignment requiring \blanknum{45}.\par
\par Well, the day before Christmas break I was told to meet with the guidance counsellor (辅导员). My mind raced with \blanknum{46}, thinking, "What could I possibly have done wrong?" But when I got there, there was a huge wrapped box on the table. The counsellor told me it was for me and \blanknum{47} me to open it. It was full of art \blanknum{48}: coloured pencils, pens, etc.\par
\par I remember \blanknum{49} and crying. I asked who got it for me and the counsellor told me "Mr. Good". I insisted on knowing who it was so I could \blanknum{50} them. She refused to tell me. Those gifts helped me create art for years to come and I graduated with the \blanknum{51} — Most Artistic Person.\par
\par I get all \blanknum{52} thinking about this every now and then. Those supplies were not just tools for drawing—they were a(n) \blanknum{53} of hope and belief in me. That simple act of kindness \blanknum{54} a seed that has grown into a lifelong \blanknum{55} for art.\par
\cz{41}{designing}{judging}{introducing}{discovering}
\cz{42}{future}{place}{privacy}{safety}
\cz{43}{suddenly}{gradually}{barely}{really}
\cz{44}{debate}{pair}{negotiate}{compare}
\cz{45}{speaking}{writing}{drawing}{reading}
\cz{46}{worry}{disappointment}{anger}{sympathy}
\cz{47}{directed}{forced}{encouraged}{ordered}
\cz{48}{commercials}{supplies}{paintings}{brochures}
\cz{49}{showing up}{hanging out}{pulling through}{breaking down}
\cz{50}{accept}{support}{recognise}{thank}
\cz{51}{assignment}{honour}{reward}{anticipation}
\cz{52}{sharp-eared}{clear-headed}{misty-eyed}{tight-lipped}
\cz{53}{track}{summary}{absence}{symbol}
\cz{54}{sowed}{lent}{moved}{wrapped}
\cz{55}{enthusiasm}{demand}{memory}{respect}
\sectitle{第二节（共10小题；每小题1.5ladnepen分，满分15分）}
\par 阅读下面短文，在空白处填入1个适当的单词或括号内单词的正确形式。\par
\par Ieoh Ming Pei, a well-known Chinese American architect, celebrated for his innovative designs and great influence on modern architecture, made significant contributions to the combination of Eastern and Western architectural styles. One architectural critic \blanknum{56} (enthusiastic) remarked, "Pei's work demonstrates how modern structures can coexist harmoniously with classical architecture, creating \blanknum{57} dialogue between the past and the present."\par
\par With the vision of combining cultural elements at its core, the Suzhou Museum, \blanknum{58} (design) through the efforts of Pei and his team, became a vivid \blanknum{59} (prove) of his architectural philosophy. The museum, with its elegant combination of traditional Chinese gardens and modernist features, \blanknum{60} (draw) over 3 million visitors in its first year after opening, many of \blanknum{61} engaged in workshops that explored the intersection of traditional Chinese art and modern design. Pei was a consistent advocate for architecture \blanknum{62} a medium of cultural exchange, emphasizing that buildings should not merely serve functional purposes \blanknum{63} act as bridges connecting diverse cultures as well. \blanknum{64} (serve) as evidence of harmony between Eastern and Western styles, Pei's legacy (遗产) highlights the potential for architecture \blanknum{65} (bridge) cultural gaps and inspire appreciation across generations as the world becomes increasingly interconnected.\par
\parttitle{第四部分 写作（共两节，满分40分）}
\sectitle{第一节（满分15分）}
\par 假定你是李华，上周参加了你校举办的“文明礼伙伴我行”(Etiquette Guides My Journey)英文手抄报设计大赛。请你给新西兰笔友Chris写一封邮件分享这次经历，内容包括：\par
\par (1) 你设计的作品;\par
\par (2) 你的理由。\par
\par (1) 写作词数应为 80 左右;\par
\par (2) 请按如下格式在答题纸的相应位置作答。\par
\par Dear Chris,\par
\writinglines{4}
\writinglines{4}
\writinglines{4}
\writinglines{4}
\par Yours, Li Hua\par
\sectitle{第二节（满分25分）}
\par 阅读下面材料，根据其内容和所给段落开头语续写两段，使之构成一篇完整的短文。\par
\par In my freshman year of college, I had a habit of sitting in the same spot every day in my economics class. It was a comfortable little corner near the window, where the sunlight would stream in and warm the desk. I loved this spot because it gave me a sense of routine and comfort during the busy college life.\par
\par One day, as I settled into my usual seat, I noticed a new face sitting next to me. She was a foreign exchange student named Mia, who had just arrived from Italy. Her English was basic, and the most complex sentence I had heard her say was, "I love this sunny weather!" Despite the language barrier, Mia always greeted me with a bright smile and a cheerful "Morning, Alex!" every day.\par
\par Mia had a habit of spreading her belongings across my desk space. Her notebooks, snacks, and even her colourful scarf would somehow find their way onto my side. Each time I walked into class, Mia would quickly realize her mistake, flash an apologetic smile, and say, "Oh, Alex! Sorry, sorry!" while hurriedly clearing my space. She would then give a nod and ask, "Ready for class, yeah?"\par
\par Initially, I found this routine mildly annoying. I thought to myself, "Why can't she keep her things on her side? Doesn't she know I sit here every day?" The last thing I wanted was to deal with a guy who barely spoke my language at eight in the morning. Just get your stuff off my desk!\par
\par One particularly busy morning, I arrived a few minutes late to class. As I stood outside the classroom, finishing a quick phone call, I noticed through the window that my desk was, as usual, occupied by Mia's things. However, before I could enter, another student, also running late, walked into the room and headed straight for my seat.\par
\par 注意：\par
\par (1) 续写词数应为 150 左右;\par
\par (2) 请按如下格式在答题卡的相应位置作答。\par
\writingpara{Paragraph 1: I watched as Mia's eyes followed the newcomer.}
\writinglines{8}
\writingpara{Paragraph 2: Then I realized Mia didn't mean to inconvenience me.}
\writinglines{8}
\end{document}
